\section{Overview and motivation}
\label{sec:overview}

Zip$^+$ is the polynomial commitment scheme used by Zinc$^+$. It commits
to a multilinear polynomial whose evaluations live in some semiring
$\mathrm{Eval}$. The prover encodes each row of the evaluation matrix
with a linear error-correcting code, arranges the encoded rows into
a matrix, and builds a Merkle tree over the columns. The two
linear-code families already shipped are:

\begin{itemize}
  \item \textbf{IPRS}: a radix-8 pseudo Reed--Solomon NTT over
    $\F_{65537}$, structurally lifted to $\Z$ so that the encoder
    operates on small integers and avoids field reductions on the
    prover hot path. This is the workhorse for the SHA / ECDSA
    integer-side lanes.
  \item \textbf{RAA}: a repeat-accumulate-accumulate code with two
    random permutations. Cheaper than IPRS but with weaker distance
    properties.
\end{itemize}

\paragraph{The motivation.}
SHA-256 and ECDSA witnesses are \emph{natively binary}: each evaluation
is most naturally represented as a polynomial of degree $<32$ over
$\F_2$ (\code{BinaryPoly<32>} in the codebase). When these witnesses
are passed through IPRS, the binary nature of the input is invisible
to the encoder; the algebra runs over $\F_{65537}$ even though every
input has only $\{0,1\}$ coefficients. A natural question is
whether a Reed--Solomon code defined directly over the binary
extension field $\GF(2^{32})$ would be a tighter fit.

\paragraph{The construction at a glance.}
Pick an irreducible $f \in \F_2[X]$ of degree $32$. Let $\flift$ be
the integer polynomial with the same monomials and coefficients in
$\{0,1\} \subset \Z$. Take a standard additive Reed--Solomon FFT
algorithm over $\GF(2^{32}) = \F_2[X]/f$ --- we use the LCH
(Lin--Chung--Han) variant with the Cantor basis, since the latter
makes the subspace polynomials $F_2$-linearised, hence
$\{0,1\}$-coefficient even after the lift. Now apply the
\emph{exact same algorithm}, with the same twiddles, but with every
arithmetic operation performed in $\Rlift$ instead of $\GF(2^{32})$.
Inputs (\code{BinaryPoly<32>}) lift trivially: each $\F_2$-coefficient
becomes a $\Z$-coefficient in $\{0,1\}$.

The resulting code has two pleasant properties:
\begin{enumerate}
  \item \textbf{Mod-2 commutativity.} Reducing the lifted codeword
    coefficient-wise modulo $2$ recovers exactly the
    $\GF(2^{32})$-RS codeword. This is the soundness witness we
    rely on to argue that the lifted code inherits good distance
    from the original RS code, modulo the lift.
  \item \textbf{Sparse-twiddle multiplication.} Every twiddle in the
    FFT is a Cantor-basis vector or a subspace-polynomial value at
    a basis vector --- all of these are lifted $\GF(2^{32})$ elements
    with $\{0,1\}$ coefficients. Multiplying any element of
    $\Rlift$ by such a twiddle is a sparse shift-and-add, no
    coefficient multiplications.
\end{enumerate}

\paragraph{Scope of this document.}
We cover (i) the algorithm, (ii) the Rust implementation in
\code{zip-plus/src/code/binary\_add\_fft/}, (iii) the
\code{ZipTypes} plumbing that makes the no-alpha-projection variant
of Zip$^+$ admissible, (iv) the protocol-level gap that currently
prevents end-to-end prove/verify, (v) the optimization journey from
naive to $\sim 26\times$-faster commit, and (vi) measured proof-size
properties (raw and zstd-3 compressed) of column openings.
